% Source: https://artofproblemsolving.com/community/c3h1669193_very_easy_problem
% Stable identity: aops-topic-1669193-q01
First check that $(8.2,0)$ lies on the line: $5\cdot8.2+12\cdot0=41$. Solving for $y$, $y=-\dfrac{5}{12}x+\dfrac{41}{12}$, so the slope is $-\dfrac{5}{12}$.

\begin{asy}
usepackage("amsmath");
size(6cm);
defaultpen(fontsize(10pt));
pair P=(8.2,0),Y=(0,41/12);
draw((-1,0)--(10.5,0),gray,Arrow(TeXHead));draw((0,-1)--(0,4.8),gray,Arrow(TeXHead));
label("$x$",(10.5,0),E,gray);label("$y$",(0,4.8),N,gray);
draw((-0.8,(41+4)/12)--(10.2,(41-51)/12),linewidth(0.9));
dot(P);
label("$(8.2,0)$",P,SW);
label("$5x+12y=41$",(6.8,2.4));
label("\mbox{}",(min(currentpicture,user=true)+max(currentpicture,user=true))/2,black); // keeps the text after the figure black when LaTeX typesets the labels inline
\end{asy}

The point-slope form through $(x_1,y_1)$ with slope $m$ is $y-y_1=m(x-x_1)$. With $(x_1,y_1)=(8.2,0)$:
\[
\boxed{y-0=-\dfrac{5}{12}(x-8.2)}.
\]
